%=============================================================================
% MÓDULO 05 — Implicações para Autoria e Responsabilidade Ética
%=============================================================================
% Padrão: As implicações devem ser fundamentadas exclusivamente nas duas fontes
% auditáveis (Clark & Chalmers, 1998 e Verbveek, 2005). Qualquer afirmação sobre
% autoria, responsabilidade ou ética deve ser rastreável a estas duas fontes.
%=============================================================================

% --- 5.1. AUTORIA ESTENDIDA HUMANO-PRÓTESE ---
% O modelo de autoria estendida alinha-se ao framework de "integração adequada" de
% Clark & Chalmers (1998) - Fonte Auditável #1 - e ao framework de "mediação
% tecnológica" de Verbveek (2005) - Fonte Auditável #2.

% A autoria não é mais individual ou do sistema isoladamente, mas estendida:
% - O humano direciona a prótese (define objetivos, restringe busca, aprova saídas)
% - A prótese executa a orquestração agente (agregação de habilidades, roteamento, execução TDD)
% - A assinatura cognitiva resultante reflete ambos, mas de modo não-redutível a nenhuma das partes

% > \textbf{Critério da Banca:} A banca verificará se o modelo de autoria estendida
% está fundamentado nas duas fontes auditáveis. A não menção a Clark & Chalmers (1998)
% e Verbveek (2005) como base para a autoria estendida será considerada inconsistência.

% --- 5.2. ACCOUNTABILITY DISTRIBUÍDO ---
% Quando a prótese falha (produz saída incorreta, viés, erro de execução), a
% responsabilidade é distribuída entre três atores, fundamentado em Verbveek (2005):

% 1. Do designer da prótese: mecanismos de gate, trust scoring, spec validation
% 2. Do usuário da prótese: supervisão adequada, definição de specs, aprovação de saídas
% 3. Do agente específico que falhou: operação fora de parâmetros esperados

% > \textbf{Critério da Banca:} O modelo de accountability tríplice deve citar explicitamente
% Verbveek (2005) como fonte fundamentada. A ausência desta referência na discussão
% da responsabilidade será bloqueador para aprovação.

% --- 5.3. APOSTAS E LIMITES ---
% O Core adota postura \textit{fail-closed}: se um executor falha ou está ausente,
% informa a falha e prossegue com recursos disponíveis. Cada falha gera reflexão
% registrada no MetaBus e ajuste de trust. Isso estabelece accountability da prótese
% em vez de culpabilização individual.

% > \textbf{Critério da Banca:} A postura fail-closed do Core e o ajuste de Trust Engine
% scores após falhas devem ser articulados com Verbveek (2005) sobre mediação tecnológica.
% A não ser por esta fonte, o modelo de accountability seria considerado desfundamentado.

% --- 5.4. LIMITES ETICOS E METODOLÓGICOS ---
% Anti-overclaim: Esta obra \textbf{não} declara que a assinatura cognitiva equivale a
% "personalidade do sistema" ou "consciência". Ela é interpretada como "rastreador de
% extensão protética", fundamentado exclusivamente nas duas fontes auditáveis.

% > \textbf{Critério da Banca:} A não ser por referência explícita às duas fontes auditáveis
% (Clark & Chalmers, 1998 e Verbveek, 2005), nenhuma declaração sobre propriedades
% mentais do sistema será aceita. Esta é uma exigência rigorosa da banca.

% --- RESUMO DAS IMPLICACÕES PARA A BANCA ---
% Todas as implicações da Seção 5 devem ser rastreáveis às duas fontes auditáveis.
% A banca verificará:
% 1. Presença de fichas críticas de Clark & Chalmers (1998) e Verbveek (2005) nas
%    seções correspondentes (Módulo 03, Seções 2.1 e 2.2)
% 2. Modelos de autoria e accountability articulados com estas duas fontes
% 3. Ausência de qualquer declaração de "mente", "personalidade" ou "consciência"
%    do Core, exceto como "rastreador de extensão protético" fundamentado nas duas fontes
% 4. Postura fail-closed do Core articulada com Verbveek (2005)